% Learning InitechBase, Chapter 2: Looking at Your Data (AN-16).  Every panel
% replayed (examples/l2-*.dot); sources cmd/query.c, core/eval.c.
\chapter{Looking at Your Data}
\Gates{test-pubs-examples, test-samir-query, test-xbase-eval}
A database answers questions only when you ask for the right fields and the
right records.  In this lesson you choose the fields \cmd{LIST} shows, list
values that InitechBase calculates, and describe the records you want with
conditions.

\begin{ChapterCovers}
\item Choosing fields
\item Comparing numbers, dates and character strings
\item Logical fields and \texttt{.NOT.}
\item Joining conditions with \texttt{.AND.} and \texttt{.OR.}
\item Listing calculated values
\item Turning the heading off
\item Memo fields
\end{ChapterCovers}

\ATEHead{Preparing for This Chapter}
Start InitechBase with a fresh copy of the sample files, as for Chapter~1, and
open the staff database:
\begin{Steps}
\item Type \type{USE STAFF} and press \key{Enter}.
\end{Steps}

\ATEHead{Choosing Fields}
After \cmd{LIST} you name the fields you want, separated by commas, in the
order you want them.  Field names are the names of the structure of the
database; Chapter~1 of \doctitle{Using InitechBase} lists those of
\file{STAFF.DBF}.
\begin{Steps}[resume]
\item Type \type{LIST LASTNAME, HIRED, SALARY} and press \key{Enter}.
\end{Steps}
\DotExample{examples/l2-fields}
Each column is as wide as its field.  Dates appear in the form month, day,
year, and salaries with their two decimal places.

\ATEHead{Comparing Numbers}
A condition compares a field with a value.  The comparison operators are
\texttt{=} equal, \texttt{<>} not equal, \texttt{<} less than, \texttt{>}
greater than, \texttt{<=} less than or equal and \texttt{>=} greater than or
equal.
\begin{Steps}[resume]
\item Type \type{LIST LASTNAME, SALARY FOR SALARY >= 40000} and press
\key{Enter}.
\ATEResponse{InitechBase lists the employees who earn 40,000 or more.}
\end{Steps}
\DotExample{examples/l2-compare}
A number in a condition is typed without commas and without a currency sign.

\ATEHead{Logical Fields}
The field \texttt{EXEMPT} is a logical field: it is true for an employee who
is exempt from overtime and false for one who is not.  A logical field is a
condition by itself.  \texttt{.NOT.} turns true into false.
\begin{Steps}[resume]
\item Type \type{LIST LASTNAME, DEPT FOR EXEMPT} and press \key{Enter}.
\item Type \type{LIST LASTNAME, DEPT FOR .NOT. EXEMPT} and press \key{Enter}.
\end{Steps}
\DotExample{examples/l2-logical}
Every employee appears in one list or the other.

\ATEHead{Comparing Dates}
A date in a condition is written as a character string in the form month, day,
year, converted with the function \fn{CTOD()}, which turns characters into a
date.  An earlier date is less than a later one.
\begin{Steps}[resume]
\item Type \type{LIST LASTNAME, HIRED FOR HIRED < CTOD('01/01/85')} and press
\key{Enter}.
\ATEResponse{InitechBase lists the employees hired before 1985.}
\end{Steps}
\DotExample{examples/l2-dates}

\ATEHead{Joining Conditions}
\texttt{.AND.} joins two conditions that must both be true; \texttt{.OR.} two
conditions of which one at least must be true.  The periods on either side
belong to the words.
\begin{Steps}[resume]
\item Type \type{LIST LASTNAME, DEPT, SALARY FOR DEPT = 'SOFT' .AND. SALARY <
40000} and press \key{Enter}.
\end{Steps}
\DotExample{examples/l2-and}
\begin{Steps}[resume]
\item Type \type{LIST LASTNAME, DEPT FOR DEPT = 'ACCT' .OR. DEPT = 'FACL'} and
press \key{Enter}.
\end{Steps}
\DotExample{examples/l2-or}
\begin{Tip}
A condition such as \texttt{DEPT = 'SOFT'} is true for every record whose
department begins with \texttt{SOFT}.  When you compare character fields,
compare with the whole value, in capital letters if the field holds capitals.
\end{Tip}

\ATEHead{Listing Calculations}
The list after \cmd{LIST} may hold calculations as well as fields.  InitechBase
works out the value for each record and uses the calculation as the heading of
its column.
\begin{Steps}[resume]
\item Type \type{LIST LASTNAME, SALARY, SALARY * 0.015} and press \key{Enter}.
\ATEResponse{InitechBase lists one and one half per cent of each salary.}
\end{Steps}
\DotExample{examples/l2-expr}
Character fields can be joined with \texttt{+}.  The function \fn{TRIM()}
removes the blanks at the end of a field, so that the first name and the last
name are separated by one blank only:
\begin{Steps}[resume]
\item Type \type{LIST TRIM(FIRSTNAME) + ' ' + LASTNAME FOR DEPT = 'ADMN'} and
press \key{Enter}.
\end{Steps}
\DotExample{examples/l2-names}

\ATEHead{Turning the Heading Off}
\begin{Steps}[resume]
\item Type \type{SET HEADING OFF} and press \key{Enter}.
\item Type \type{LIST LASTNAME, EXT FOR DEPT = 'CUST'} and press \key{Enter}.
\end{Steps}
\DotExample{examples/l2-heading}
SET HEADING ON brings the heading back.  It is also restored when you leave
InitechBase and start it again.

\ATEHead{Memo Fields}
The field \texttt{NOTES} is a memo field.  Its text is kept in the memo file
\file{STAFF.DBT}, which \cmd{USE} opened with the database.  \cmd{LIST} and
\cmd{DISPLAY} show the text.
\begin{Steps}[resume]
\item Type \type{GO 12} and press \key{Enter}.
\item Type \type{DISPLAY NOTES} and press \key{Enter}.
\end{Steps}
\DotExample{examples/l2-memo}

\ATEHead{Where Do You Go From Here?}
You can now list any fields of the records that meet a condition.  In
Chapter~3 you learn how InitechBase keeps its place in a database and how you
move from record to record.

\ATEHead{Summary}
\begin{enumerate}
\item \cmd{LIST} \xp{field list} lists the fields named, in the order named.
\item \kw{FOR} \xp{condition} lists only the records for which the condition is
  true.
\item The comparisons are \texttt{=}, \texttt{<>}, \texttt{<}, \texttt{>},
  \texttt{<=} and \texttt{>=}; \fn{CTOD()} writes a date in a condition.
\item A logical field is a condition by itself; \texttt{.NOT.}, \texttt{.AND.}
  and \texttt{.OR.} combine conditions.
\item The list may hold calculations; \texttt{+} joins character values;
  \fn{TRIM()} removes trailing blanks.
\item SET HEADING OFF leaves out the heading.
\end{enumerate}
